%% ==========================================================================
\section{What kind of claim this is}
\label{sec:method}
%% ==========================================================================

\begin{remark}[What kind of claim these are --- a praxeological reading]\label{rem:praxeology}
Neither the Race nor the Spiral is an empirical hypothesis awaiting corroboration, and
it is worth saying plainly why, because the alternative reading invites a demand this
domain cannot satisfy.

Both are deduced from the fact of purposive action under a stated adversary: \emph{given}
$\mathrm{WDY}[W,c]$, these are the conditions under which coercion is prolonged or the
holder is removed. That is the posture Mises named \emph{praxeology} --- theorems derived
from action rather than induced from observation, with history able to illustrate a
theorem but never to verify or falsify it~\cite{mises-human-action}. We claim the
posture, not the whole method. The decision rule here is quantitative, and a strict
praxeologist would resist it: $p_A$ is what Mises called a \emph{case} probability --- the
likelihood of this attacker prevailing in this encounter --- and he denied that numbers
belong there. The concession costs little, because nothing in the argument requires
$p_A$ to be \emph{computable}; only that the ordering it expresses exists. What carries
over is the epistemology, and that is the part that governs how the results should be
read.

No observation could settle either claim, and the reason is internal to the subject. The
only witnesses to an enacted Deadly Race are an attacker with every reason to deny the
motive and a holder --- or, where the Lemma's prediction is realized, a family --- with
every reason to deny the asset. Asking a convicted murderer whether he killed to win a
race for a stash whose existence the survivors also deny is not a weak instrument but a
void one: both answers are the negative that NNPP (Assumption~\ref{lem:nnpl}) holds
cannot be disproved.

The Spiral fares no better under the same test. Ask a convicted kidnapper whether the
coercion he prolonged did in fact reach a stash whose existence his victim still denies:
an affirmative confesses a robbery on top of the abduction, and a \emph{truthful} denial
is indistinguishable from the answer a successful robber would give --- while the
victim's own denial, which is what would corroborate either, is once more the negative
that cannot be disproved. \S\ref{par:ratirl} works that fork through on the one case
where it can be posed concretely: three readings stay simultaneously plausible, and each
indicts obscurity. Neither witness can be believed, and no third party was present.
The instrument that would test the model is disabled by the very mechanism the model
describes.

The standard that does apply is \emph{conditional soundness}, and it is checkable
without data. The theory can be shown wrong, in two specific ways. \textbf{Break the
derivation:} exhibit a reachable state that is sufficient for $\Kset$ but unfinished,
with no access-dominating delegate, in which $\pi = 0$. \textbf{Break the \emph{a
fortiori} step:} show that the modeled adversary is, in some respect that bears on the
result, \emph{weaker} than a real coercer --- for a bar cleared against an adversary
granted more than reality gives him (Assumption~\ref{ass:power}) is cleared against the real one
only while that grant holds. Neither route requires a case, a victim or a survey; both
are settled by reading the argument.

What is genuinely undecidable is only the per-case attribution --- whether \emph{this}
killing was enacted for \emph{that} payoff --- and the paper never asserts it. This is
also why the output is a \textbf{bar} rather than a finding. Whether a protocol leaves a
sufficient-but-unfinished state is decided by reading the protocol; whether a norm
depletes denial credibility is decided by reading the counsel. Neither requires a
victim, and that is the point of the whole construction.
\end{remark}
